\section{PRELIMINARY DISCUSSION}

\subsection{Interpretation of the Results}

Across all the simulated systems, Explicit Euler consistently showed the poorest conservation of both energy and
angular momentum, meanwhile Leapfrog, Velocity Verlet and Symplectic Euler demonstrated substantially greater
long-term stability.
Although RK4 produced the smallest energy errors at sufficiently small timesteps, its behaviour deteriorated
rapidly at coarse resolutions, highlighting the importance of timestep selection for higher-order
non-symplectic methods.

Explicit Euler demonstates such behaviour due to the accumulated truncation errors over each timestep.
Symplectic methods by their mathematical design preserve the Hamiltonian geometry.
Instead of estimating a precise value of energy, symplectic methods revolve around the true value in the sinusoid.
For example, the RK4 algorithm minimises the local truncation error with the higher resolution, yet doesn't preserve
the phase space.
Hence why, there is a difference between local accuracy and long-term stability.

\subsection{Relation to Hamiltonian Mechanics}

The observed behaviour closely matches theoretical expectations for Hamiltonian systems.
Symplectic methods do not conserve energy exactly; instead they preserve a nearby modified Hamiltonian,
resulting in bounded oscillatory error rather than secular drift.
This distinction is central to backward error analysis: a symplectic integrator can be shown to solve, up
to exponentially small corrections in the timestep, the flow of a shadow Hamiltonian $\tilde{H}$ that lies close
to the true Hamiltonian $H$. Because $\tilde{H}$ is itself exactly conserved by the discrete map, the numerically
computed energy oscillates within a fixed band about the true value rather than drifting away from it.
Non-symplectic methods such as RK4 possess no equivalent conserved quantity, so their local truncation errors
are free to accumulate coherently over long integrations, producing the secular drift observed at coarse
resolutions.

\subsection{Angular Momentum Conservation}

The conservation of angular momentum provides a complementary, and in some respects more discriminating, test
of integrator quality than energy alone. In the continuous two-body problem, angular momentum
$\mathbf{L} = \mathbf{r} \times \mathbf{p}$ is conserved exactly because the gravitational potential depends
only on the separation $|\mathbf{r}|$ and is therefore invariant under rotations; by Noether's theorem, this
rotational symmetry guarantees a corresponding conserved quantity.

The discrete integrators inherit this symmetry to varying degrees. Symplectic Euler, Leapfrog and Velocity
Verlet update velocity using a force evaluated at the current (or half-step) position, and this force is
always directed along the instantaneous radius vector. Because the discrete torque $\mathbf{r} \times
\mathbf{F}(\mathbf{r})$ vanishes identically at every timestep, exactly as it does in the continuous system,
these methods conserve angular momentum to machine precision regardless of the timestep chosen. This is a
structural property of the update rule rather than an accuracy effect, and it explains why angular momentum
conservation for the symplectic methods did not degrade even where energy error grew at coarse resolutions.
This structural distinction was formally established by \citet{zhang1995}, who showed that symplectic
integrators conserve angular momentum exactly while nonsymplectic integrators generally do not; the results
obtained here for the two-body problem are consistent with that finding.

Explicit Euler and RK4 do not share this property. Explicit Euler updates position and velocity from
inconsistent, staggered information (the velocity used to advance position is not the velocity implied by
the force just computed), which introduces a small effective torque at each step; this error compounds
secularly, matching the drift observed in energy. RK4, despite averaging force evaluations across several
sub-stages within a step, does not evaluate the force along a single self-consistent trajectory, so it also
fails to preserve $\mathbf{r} \times \mathbf{F} = \mathbf{0}$ exactly at the discrete level. Its angular
momentum error was nonetheless substantially smaller than that of Explicit Euler at fine timesteps, consistent
with its higher local order of accuracy, but the error was not bounded in the same structural sense as the
symplectic methods and would be expected to grow if the timestep were coarsened further or the integration
extended over longer timescales. Taken together, the angular momentum results reinforce the central conclusion
of this study: symplectic structure, not local order of accuracy, is the primary determinant of long-term
qualitative fidelity in Hamiltonian simulations.

\subsection{RK4 Behaviour}

RK4 possesses fourth-order local accuracy, but this accuracy assumes sufficiently small timesteps.
At coarse resolutions, the large integration interval dominates truncation error and the method rapidly
loses stability. Once the timestep becomes sufficiently small, the fourth-order convergence dominates,
allowing RK4 to outperform all remaining methods in terms of instantaneous energy error.

However, despite this excellent short-term accuracy, RK4 remains non-symplectic and therefore cannot guarantee
long-term conservation of invariants in Hamiltonian systems.

\subsection{Computational Cost v. Accuracy}

Although RK4 produced the lowest numerical error at fine resolution, this came at approximately four force evaluations
per timestep and the highest runtime among all methods.
Leapfrog and Velocity Verlet therefore provide a considerably better compromise between computational efficiency and
conservation accuracy.

\subsection{Velocity Verlet v. Leapfrog}

Velocity Verlet and Leapfrog are algebraically equivalent formulations of the same second-order symplectic
method, differing only in how the half-step velocity is stored and reported, and the results confirm this
equivalence at every level examined. In runtime, Leapfrog was marginally the fastest integrator at higher
resolutions, with Velocity Verlet close behind it; the two are not expected to differ by more than
implementation overhead, since both require one force evaluation per timestep. In energy conservation,
Leapfrog was found to overlap almost exactly with Velocity Verlet on the error-versus-timestep comparison,
with both methods consistently outperforming Symplectic Euler by roughly two orders of magnitude at a given
resolution. The angular momentum results show the same equivalence in its clearest form: both integrators
produced an identical mean absolute angular momentum value across all three tested resolutions, and their
relative angular momentum error sat at the floating-point noise floor of approximately $10^{-13}$,
independent of timestep. This lends strong empirical support to the theoretical expectation that the two
schemes trace out the same discrete trajectory, and any residual difference between them observed here is
attributable to implementation detail (such as the point in the step at which velocity is reported) rather
than to a genuine difference in the underlying numerical method. Consequently, neither integrator can be
said to be superior to the other on the criteria examined in this study; the choice between them in practice
would come down to convenience, for instance whether a synchronised position-velocity pair at each full step
(as reported by Velocity Verlet) is more useful downstream than the staggered representation used internally
by Leapfrog.

\subsection{Hypothesis}

The hypothesis was only partially supported. Velocity Verlet demonstrated an excellent balance between runtime
and conservation accuracy, closely matching Leapfrog due to their mathematical equivalence.
However, RK4 produced lower energy errors at sufficiently fine timesteps, although at substantially greater
computational cost and without preserving the symplectic structure.

\subsection{Limitations}

While the results presented here are consistent with established theory on symplectic integration, several
limitations constrain the scope of the conclusions and point towards natural extensions of this work.

\paragraph{Fixed timestep integration.} All integrators were evaluated at a single, fixed timestep rather than
with an adaptive scheme. A fixed $\Delta t$ is necessarily a compromise: it must be small enough to resolve
the fastest dynamics of the orbit, typically near perihelion where the angular velocity is greatest, while
remaining computationally tractable over the full integration window. An adaptive-timestep approach, whether
based on local truncation error estimation (as in embedded Runge--Kutta pairs) or on a time-regularising
coordinate transformation such as the Sundman transformation, would allow the timestep to shrink automatically
in regions of high curvature and expand elsewhere, and would likely narrow the performance gap observed
between RK4 and the symplectic methods at coarse resolutions. Adapting a symplectic method to variable
timesteps is non-trivial, however, since naive step-size control generally destroys the symplectic property
of the map, and this trade-off was not explored here.

\paragraph{Higher-order symplectic compositions.} The symplectic methods considered in this study (Symplectic
Euler, Leapfrog and Velocity Verlet) are all first- or second-order accurate. Higher-order symplectic
integrators can be constructed by composing second-order symplectic maps with carefully chosen step-size
coefficients, as in the fourth-order schemes of Forest and Ruth or Yoshida. Such methods would combine the
long-term structural stability demonstrated here with markedly reduced local error, and a comparison against
them would clarify whether the accuracy advantage of RK4 at fine timesteps persists against a symplectic method
of matching order.

\paragraph{Restriction to a two-body system.} The analysis was restricted to an isolated two-body Hamiltonian
system, for which an analytical Keplerian solution exists and can be used as a ground truth. Extending the
comparison to $N$-body simulations, where no closed-form solution is generally available, would test whether
the relative ranking of integrators established here persists in the presence of additional gravitational
interactions, and would allow investigation of numerical performance in genuinely chaotic orbital
configurations, where nearby trajectories diverge exponentially and long-term structural conservation becomes
even more important than local accuracy. Related settings of particular interest include planetary systems
exhibiting mean-motion resonances, where small integration errors can be resonantly amplified over long
timescales, and the long-term stability analysis of exoplanetary systems, which typically requires integration
over timescales of $10^{6}$--$10^{9}$ orbits and is correspondingly sensitive to the choice of integrator.

\paragraph{Computational scaling.} The present implementation was run in serial on three cores of a consumer device,
and no attempt was made to characterise performance at scale. $N$-body extensions of this work would benefit from
GPU acceleration, which parallelises the pairwise force calculations that dominate the computational cost of
large simulations, and from shared- or distributed-memory parallelisation using OpenMP or MPI respectively.
A systematic scaling study would be needed to establish whether the computational-cost conclusions drawn here
for a single two-body pair generalise to the substantially larger force-evaluation counts required by $N$-body
and resonance studies.

\paragraph{Newtonian approximation.} All simulations assumed Newtonian gravity, neglecting relativistic
corrections. For systems with strong or close-orbit gravitational fields, such as tight binary star systems
or bodies passing close to a compact object, post-Newtonian corrections introduce additional effects,
including apsidal precession, that a purely Newtonian integrator cannot reproduce regardless of its symplectic
properties. Incorporating these corrections would require modifying the force law itself rather than the
integration scheme, and represents a natural extension once the present comparison is established on the
simpler Newtonian case.